\begin{solapartat}
\punts{0,5} Dibuix: el camp elèctric està dirigit cap a càrregues negatives i radial per la geometria esfèrica:

\figura[0.25]{sol_fig1.png}

\punts{0,5} L'enunciat ens diu que el camp elèctric és el generat per tota la càrrega al centre de l'esfera. Sabem que el mòdul de camp elèctric en aquest cas és:
\[ \bigl|\vec{E}\bigr| = k\,\frac{q}{r^2} \]

\punts{0,15} Per tant, la càrrega és: \quad $|q| = \frac{|E|r^2}{k} = \frac{150\cdot(6,37\ 10^{6})^2}{8,99\ 10^{9}} = 6,77\cdot 10^{5}\ \mathrm{C}$

\punts{0,10} Per considerar que la càrrega és negativa: $q = -6,77\cdot 10^{5}\ \mathrm{C}$
\end{solapartat}

\begin{solapartat}
\punts{0,25} El seu mòdul serà: $\bigl|\vec{F}_e\bigr| = q|E| = 1,602\ 10^{-19}\cdot 150 = 2,403\ 10^{-17}\ \mathrm{N}$

\punts{0,5} El pes de la gota d'aigua ha de ser igual a la força elèctrica creada pel camp ja calculada:
\[ \bigl|\vec{F}_e\bigr| = \bigl|\vec{P}\bigr| = mg \text{ i, per tant, } m = \frac{\bigl|\vec{F}_e\bigr|}{g} = \frac{2,403\ 10^{-17}}{9,81} = 2,45\ 10^{-18}\ \mathrm{kg} \]

\punts{0,25} Esquema: la força elèctrica, contrària al camp elèctric. S'oposa al pes.

\figura[0.3]{sol_fig2.png}

\punts{0,25} Per calcular el diàmetre de la gota, sabem que la massa de la gota és el seu volum per la densitat de l'aigua $m = \rho\frac{4}{3}\pi r^3$ i trobem el radi:
\[ r = \sqrt[3]{\frac{3m}{4\pi\rho}} = \sqrt[3]{\frac{3\cdot 2,45\ 10^{-18}}{4\pi\cdot 10^{3}}} = 0,836\cdot 10^{-7}\ \mathrm{m}. \]
El diàmetre és dues vegades el radi,
\[ d = 1,67\cdot 10^{-7}\ \mathrm{m} = 0,167\ \mathrm{\mu m} = 167\ \mathrm{nm} \]
\end{solapartat}
